\documentclass[10pt,a4paper]{amsart}
\usepackage[margin=19mm]{geometry}
\usepackage{amsmath,amssymb,mathtools}
\usepackage[scr=rsfs,cal=euler]{mathalfa}
\usepackage{xfrac}
\usepackage[unicode,hidelinks,bookmarks=false]{hyperref}
\newcommand{\ie}{i.e.\ }
\newcommand{\cf}{cf.\ }
\newcommand{\resp}{respectively\ }
\newcommand{\Subsection}{\subsection}
\providecommand{\nobreakdash}{\nobreak\hbox{-}}
\setlength{\emergencystretch}{2em}
\multlinegap0pt
\usepackage{bm}
\usepackage{bbm}
\usepackage{dsfont}
\usepackage{xspace}
\usepackage{cancel}
\usepackage{mathtools}
\usepackage{amsthm}
\makeatletter
\@ifundefined{lemma}{\newtheorem{lemma}{Lemma}}{}
\makeatother
\newcommand{\RR}{\mathbb{R}}
\title[Uniform positivity of Lyapunov exponents for anti H\"older p]{Uniform positivity of Lyapunov exponents for anti H\"older potentials}
\author{Nicholas Chiem}
\date{Source: arXiv:2606.10033v1 (CC BY 4.0)}
\begin{document}
\maketitle

\noindent\emph{Source and licence.} Uniform positivity of Lyapunov exponents for anti H\"older potentials. Version arXiv:2606.10033v1 (CC BY 4.0), distributed under the Creative Commons Attribution 4.0 International licence, as recorded in the arXiv licence metadata for this exact version and shown on the abstract page https://arxiv.org/abs/2606.10033v1. Full rights notice: licenses/arxiv-2606.10033-CC-BY-4.0.txt.\par\smallskip

\noindent\textit{Editorial scope.} Counted unit: the Lemma whose source label is \texttt{l.hyperplane} in the article (source0.tex lines 1151--1156, source label \texttt{l.hyperplane}), delivered together with the article's own complete original proof of it (source0.tex lines 1157--1202, one \texttt{proof} environment of 46 lines). No mathematics is written, repaired or re-derived: statement and proof are the article's own text line for line. Required context retained uncounted: none. Every definition, display and label of those lines is retained unchanged. Omitted from this package and not redistributed: 1--1150, 1203--1493. Nothing inside a retained excerpt is omitted or altered: every retained body is delivered exactly as the article's lines are written, so no omission receipt of any kind accompanies this package. Citations: the retained excerpts carry no citation command at all, so no bibliography entry of the article is transcribed and no reference is redistributed. Reference adaptation: none was needed and none was made; every reference inside the delivered unit resolves inside it, so no unresolved reference or citation is produced. Printed slips kept literal: no printed word, symbol, hypothesis or number of the counted statement or of its proof was altered. Layout and class: the article's own class is \texttt{amsart} with packages including graphicx, amsmath, amssymb, amsthm, tikz, esint, hyperref, stmaryrd, dirtytalk, pgfplots, float; the wrapper is the lane's pinned \texttt{amsart} editorial wrapper at 10pt on A4, which restates the theorem-like environments the retained text needs and adds the article's own macro definitions that the retained text needs (\texttt{\textbackslash RR}), with extra wrapper packages (bm, bbm, dsfont, xspace, cancel, mathtools). The delivered numbering is the unit's own. Assets: no retained span contains a figure environment, a graphics-inclusion command, a PDF-page inclusion, a table or a TikZ picture; no image file is copied into this package, the package declares no asset dependency and no image file is redistributed with it. Licence: version arXiv:2606.10033v1 of this article is distributed under the Creative Commons Attribution 4.0 International licence, as recorded in the arXiv licence metadata for this exact version and shown on the abstract page https://arxiv.org/abs/2606.10033v1. A full rights notice is delivered at licenses/arxiv-2606.10033-CC-BY-4.0.txt. Characters: every retained excerpt is pure ASCII, so the delivered engine, XeLaTeX with its default Latin Modern text font, reports no missing character.
\begin{lemma}\label{l.hyperplane}
    Consider a positive integer $d\geq 2$ and assume $u\in\RR^d$ is a unit vector. Then there exists an affine hyperplane $P\subset \RR^d$ such that when we define
    \[Z = P \cap F,\]
    the following statement holds: For every $\omega\in F$, there exists a unique $z\in Z$ such that $\omega\in $Image$(f_z)\cap F$ where
    \[f_z(t) = tu + z.\]
\end{lemma}

\begin{proof}
    Without loss of generality, we may assume
    \[u_i = \begin{cases}
        > 0 & i = 1,\dots, k\\
        < 0 & i = k+1, \dots, j\\
        = 0 & i = j+1, \dots, d
    \end{cases}\]
    for $0\leq k,j\leq d$ where $0\leq k + j \leq d$. Note one can permute the $1,\dots, d$ subscripts to match the original vector. 

    Define an affine hyperplane 
    \[P:= \{x\in \RR^d: \sum_{i=1}^k x_i - \sum_{i=k+1}^{k+j} x_i = 1\}.\]
    Let $Z = P\cap F$. We claim this choice of $Z$ suffices. Let $\omega\in F$ and define
    \[g(\omega)= \sum_{i=1}^k \omega_i - \sum_{i=k+1}^{k+j} \omega_i.\]
    Notice that $g(u) \neq 0$, since if $g(u) = 0$, then 
    \[\sum_{i=1}^k \omega_i = \sum_{i=k+1}^{k+j} \omega_i.\]
    The left hand side is strictly positive and the right is strictly negative, which is a contradiction. 

    Define $t_0$ as
    \[t_0 = \frac{g(\omega) - 1}{g(u)}\]
    and $z \in \RR^d$ with components defined as
    \[z_i = \begin{cases}
        \omega_i -t_0u_i & i = 1,\dots, k+j\\
        0 & i = k+j + 1, \dots, d
    \end{cases}.\]
    Observe
    \begin{align*}
        g(z) &= \sum_{i=1}^k z_i - \sum_{i=k+1}^{k+j} z_i\\
        &= \sum_{i=1}^k (\omega_i -tu_i) - \sum_{i=k+1}^{k+j} (\omega_i -tu_i)\\
        &= g(\omega) - tg(u)\\
        &= g(\omega) - (g(\omega)-1)\\
        &= 1,
    \end{align*}
    which implies that $z\in P$. 
    
    It remains to show that $z\in F$, so consider three cases:
    \begin{enumerate}
        \item $\omega\in P$,
        \item $\omega$ satisfies $g(\omega)>1$,
        \item $\omega$ satisfies $g(\omega)<1$.
    \end{enumerate}
    Case 1 is trivial as $z = \omega$. Let us consider case $2$. Assume for a contradiction that there exists a $j = 1,\dots, k$ such that $z_j < 0$. Define $\Tilde{t} = \frac{\omega_j}{u_j}$. Since $z_j < 0$,
    \[z_j = \omega_j - t_0 u_j < 0 \implies \frac{\omega_j}{u_j} < t_0 \implies \Tilde{t} < t_0.\] 
    But $\omega_j - tu_j$ is a decreasing function and $w_j -\Tilde{t}u_j = 0$, so $\Tilde{t}> t_0$, which is a contradiction.
    
    The argument for $i = k+1, \dots, k+j$ is the same as above and there is nothing to show when $i = k+j+1, \dots, d$. Also, notice if $z\in Z$, then $f_z(t_0) = \omega$ and hence $\omega\in$ Image$(f_z)\cap F$. The proof for case 3 has an analogous argument to case 2.
\end{proof}

\end{document}
